\documentclass[11pt,a4paper]{article}
\usepackage[utf8]{inputenc}
\usepackage{amsmath,amssymb,amsthm}
\usepackage{algorithm}
\usepackage{algpseudocode}
\usepackage{booktabs}
\usepackage{hyperref}
\usepackage{geometry}
\geometry{margin=1in}

\title{Pointer Networks and Copy Mechanisms: Direct Positional Indexing and Hybrid Generative-Extractive Soft Switching}
\author{Team Scaibu \\ \small Email: \texttt{jaydeep.9664920749@gmail.com} \\ \small Architecture and Core Engine Team}
\date{\today}

\newtheorem{theorem}{Theorem}
\newtheorem{lemma}[theorem]{Lemma}
\newtheorem{proposition}[theorem]{Proposition}
\theoremstyle{definition}
\newtheorem{definition}{Definition}
\newtheorem{remark}{Remark}

\begin{document}

\maketitle

\begin{abstract}
Standard sequence decoders generate tokens strictly from a pre-determined, fixed-size vocabulary $\mathcal{V}_{\text{fixed}}$, rendering them incapable of emitting novel named entities, numerical constants, or solving variable-length combinatorial indexing tasks (such as Convex Hull, Delaunay Triangulation, and Traveling Salesperson problems). Pointer Networks resolve this limitation by replacing fixed-vocabulary softmax projections with direct attention pointing distributions over input sequence positions. Pointer-Generator architectures generalize this mechanism into a hybrid framework via a continuous generation gate $p_{\text{gen}} \in [0, 1]$, smoothly interpolating between abstract vocabulary generation and verbatim source token copying. This document formalizes both mechanisms, proves probability conservation across extended vocabularies, details algorithmic pseudocode, and delivers a verified numerical calculation.
\end{abstract}

\section{Notation and Mathematical Preliminaries}

Let $\mathbf{x} = (x_1, \dots, x_{T_x})$ denote the source token sequence where $x_i \in \mathcal{V}_{\text{ext}}$, $\mathbf{s}_t \in \mathbb{R}^{d_s}$ denote the decoder state at step $t$, and $\mathbf{c}_t \in \mathbb{R}^{d_c}$ denote the attention context vector.

\begin{table}[h]
\centering
\caption{Mathematical Notation for Pointer-Generator Copy Networks}
\begin{tabular}{lll}
\toprule
\textbf{Symbol} & \textbf{Domain} & \textbf{Semantic Description} \\
\midrule
$\mathcal{V}_{\text{fixed}}$ & Set & Predefined fixed language vocabulary ($|\mathcal{V}_{\text{fixed}}| = V$) \\
$\mathcal{V}_{\text{ext}}$ & Set & Extended dynamic vocabulary: $\mathcal{V}_{\text{fixed}} \cup \{x_1, \dots, x_{T_x}\}$ \\
$p_{\text{gen}}$ & $[0, 1]$ & Continuous scalar probability of generating from $\mathcal{V}_{\text{fixed}}$ \\
$1 - p_{\text{gen}}$ & $[0, 1]$ & Probability of copying a token directly from the input sequence \\
$\boldsymbol{\alpha}_t$ & $\Delta^{T_x - 1}$ & Attention alignment probability vector over source positions \\
$P_{\text{vocab}}(w)$ & $[0, 1]$ & Baseline fixed-vocabulary output probability distribution \\
$P(w)$ & $[0, 1]$ & Final composite probability distribution over extended vocabulary \\
\bottomrule
\end{tabular}
\end{table}

\section{Problem Definition and Core Concepts}

\subsection{3-Level Explanation}
\begin{enumerate}
    \item \textbf{Intuitive (High School level):} When summarizing an article about an unfamiliar person (like ``Dr. Xylokopoulos''), a normal AI fails because that strange name isn't in its dictionary. A copy-pointer mechanism gives the AI a photocopier: it can either write standard English words or directly point to and photocopy the exact name from the source text.
    \item \textbf{Engineering Level:} Pointer Networks use attention weights $\alpha_t$ directly as output probabilities over input indices. Pointer-Generators compute a dynamic soft switch $p_{\text{gen}} = \sigma(\mathbf{w}_c^T \mathbf{c}_t + \mathbf{w}_s^T \mathbf{s}_t + \mathbf{w}_x^T \mathbf{x}_t + b)$, producing composite distribution $P(w) = p_{\text{gen}} P_{\text{vocab}}(w) + (1 - p_{\text{gen}}) \sum_{i: x_i = w} \alpha_{t, i}$.
    \item \textbf{Mathematical Level:} Given extended vocabulary $\mathcal{V}_{\text{ext}}$, the composite density satisfies:
    \begin{equation}
    P(w) = p_{\text{gen}} \sum_{v=1}^V \mathbb{I}(w = v) P_{\text{vocab}}(v) + (1 - p_{\text{gen}}) \sum_{i=1}^{T_x} \mathbb{I}(w = x_i) \alpha_{t, i}
    \end{equation}
    where $p_{\text{gen}} = \sigma(\mathbf{w}_c^T \mathbf{c}_t + \mathbf{w}_s^T \mathbf{s}_t + \mathbf{w}_x^T \mathbf{x}_t + b_{\text{gen}})$.
\end{enumerate}

\subsection{5-Step Algorithmic Framework}
\begin{enumerate}
    \item \textbf{Attention and Context Computation:} Evaluate attention energy scores $e_{t, i}$ and form context $\mathbf{c}_t = \sum_i \alpha_{t, i} \mathbf{h}_i$.
    \item \textbf{Vocabulary Distribution Projection:} Compute standard vocabulary logits and softmax $P_{\text{vocab}} = \text{Softmax}(W_v \mathbf{s}_t + \mathbf{b}_v)$.
    \item \textbf{Generation Probability Gating:} Calculate soft scalar gate $p_{\text{gen}} = \sigma(\mathbf{w}_c^T \mathbf{c}_t + \mathbf{w}_s^T \mathbf{s}_t + \mathbf{w}_x^T \mathbf{x}_t + b_{\text{gen}})$.
    \item \textbf{Extended Vocabulary Mixture:} Weight $P_{\text{vocab}}$ by $p_{\text{gen}}$ and accumulate attention mass $(1 - p_{\text{gen}}) \alpha_{t, i}$ at token index $x_i \in \mathcal{V}_{\text{ext}}$.
    \item \textbf{Loss Backpropagation / Output Emission:} Evaluate negative log-likelihood $\mathcal{L} = -\log P(y_t^*)$ across the extended vocabulary.
\end{enumerate}

\section{Algorithm Specification}

\begin{algorithm}[H]
\caption{Pointer-Generator Hybrid Copy Step}
\label{alg:pointer_copy}
\begin{algorithmic}[1]
\Require Decoder state $\mathbf{s}_t$, context $\mathbf{c}_t$, decoder input $\mathbf{x}_t$, fixed vocab probs $P_{\text{vocab}}$, attention weights $\boldsymbol{\alpha}_t$, source token IDs $(x_1, \dots, x_{T_x})$, extended size $V_{\text{ext}}$.
\Ensure Composite distribution $P \in \mathbb{R}^{V_{\text{ext}}}$, generation probability $p_{\text{gen}}$.
\State $p_{\text{gen}} \gets \sigma(\mathbf{w}_c^T \mathbf{c}_t + \mathbf{w}_s^T \mathbf{s}_t + \mathbf{w}_x^T \mathbf{x}_t + b_{\text{gen}})$
\State Initialize $P[w] \gets 0.0$ for all $w \in \{1, \dots, V_{\text{ext}}\}$
\For{$v = 1$ \textbf{to} $V_{\text{fixed}}$}
    \State $P[v] \gets P[v] + p_{\text{gen}} \cdot P_{\text{vocab}}[v]$
\EndFor
\For{$i = 1$ \textbf{to} $T_x$}
    \State $w \gets x_i$
    \State $P[w] \gets P[w] + (1.0 - p_{\text{gen}}) \cdot \alpha_t[i]$
\EndFor
\State \Return $P, p_{\text{gen}}$
\end{algorithmic}
\end{algorithm}

\section{Theoretical Guarantees and Probability Conservation}

\begin{theorem}[Conservation of Extended Probability Measure]
\label{thm:copy_measure}
For any $p_{\text{gen}} \in [0, 1]$, valid vocabulary distribution $\sum_{v=1}^{V_{\text{fixed}}} P_{\text{vocab}}(v) = 1$, and valid attention distribution $\sum_{i=1}^{T_x} \alpha_{t, i} = 1$, the composite hybrid measure sums identically to unity:
\begin{equation}
\sum_{w \in \mathcal{V}_{\text{ext}}} P(w) = 1.0
\end{equation}
\end{theorem}

\begin{proof}
Summing over the disjoint extended support:
\begin{align}
\sum_{w \in \mathcal{V}_{\text{ext}}} P(w) &= \sum_{w \in \mathcal{V}_{\text{ext}}} \left[ p_{\text{gen}} \sum_{v=1}^{V_{\text{fixed}}} \mathbb{I}(w = v) P_{\text{vocab}}(v) + (1 - p_{\text{gen}}) \sum_{i=1}^{T_x} \mathbb{I}(w = x_i) \alpha_{t, i} \right] \\
&= p_{\text{gen}} \sum_{v=1}^{V_{\text{fixed}}} P_{\text{vocab}}(v) + (1 - p_{\text{gen}}) \sum_{i=1}^{T_x} \alpha_{t, i} \\
&= p_{\text{gen}}(1) + (1 - p_{\text{gen}})(1) = 1
\end{align}
\end{proof}

\section{Complexity Analysis}

\begin{itemize}
    \item \textbf{Time Complexity:} $\mathcal{O}(V_{\text{fixed}} + T_x)$ linear scatter operations per decoding step.
    \item \textbf{Space Complexity:} $\mathcal{O}(V_{\text{ext}})$ memory to store the extended probability vector.
\end{itemize}

\section{Exactness and Error Analysis}

The copy mechanism is end-to-end differentiable. When a target token $y^*$ appears both in $\mathcal{V}_{\text{fixed}}$ and the source text ($y^* = x_i$), the loss gradient $\frac{\partial \mathcal{L}}{\partial p_{\text{gen}}}$ naturally coordinates whether to generate or copy based on prediction confidence.

\section{Worked Numerical Example}

Let $V_{\text{fixed}} = 2$ ($0=\text{the}, 1=\text{cat}$), $V_{\text{ext}} = 3$ ($2=\text{Aristotle}$).
Source tokens: $[x_1, x_2] = [0, 2]$ (``the Aristotle'').
Let $P_{\text{vocab}} = [0.8, 0.2]$, attention weights $\boldsymbol{\alpha} = [0.1, 0.9]$.
Let computed $p_{\text{gen}} = 0.6 \implies 1 - p_{\text{gen}} = 0.4$.

\textbf{Final Extended Probabilities:}
\begin{itemize}
    \item $w=0$ (``the''): $0.6 \times 0.8 + 0.4 \times 0.1 = 0.48 + 0.04 = 0.52$.
    \item $w=1$ (``cat''): $0.6 \times 0.2 + 0.4 \times 0.0 = 0.12$.
    \item $w=2$ (``Aristotle''): $0.6 \times 0.0 + 0.4 \times 0.9 = 0.36$.
    \item Check sum: $0.52 + 0.12 + 0.36 = 1.00$.
\end{itemize}

\section{Numerical Considerations and Edge Cases}

\begin{itemize}
    \item \textbf{Coverage Loss:} Pointer-generator summarization models frequently suffer from word repetition (copying the same phrase repeatedly). Adding a coverage loss $\mathcal{L}_{\text{cov}} = \sum_t \sum_i \min(\alpha_{t, i}, \sum_{u=1}^{t-1} \alpha_{u, i})$ penalizes repeated attention to the same source tokens.
\end{itemize}

\section{References}
\begin{enumerate}
    \item Vinyals, O., Fortunato, M., & Jaitly, N. (2015). Pointer networks. \emph{NeurIPS}, 2692--2700.
    \item See, A., Liu, P. J., & Manning, C. D. (2017). Get to the point: Summarization with pointer-generator networks. \emph{ACL}, 1073--1083.
\end{enumerate}

\end{document}
